% !TEX root = ../main.tex
\chapter{\cmm{}: tipos, operadores e controle}
\label{cap:cmm-tipos-operadores}

\section{Tipos de dados}

A \cmm{} tem exatamente quatro palavras-chave de tipo. O \code{void} marca a ausência de tipo, no retorno de funções que não retornam nada. O \code{int} é o inteiro de \code{\#NUBITS} bits, em complemento de dois. O \code{float} é o ponto flutuante no formato próprio do SAPHO, com mantissa de \code{\#NBMANT} bits e expoente de \code{\#NBEXPO} bits (Seção~\ref{sec:float-sapho}). E o \code{comp} é o número complexo, um par de \code{float} com as partes real e imaginária, tratado como tipo de primeira classe da linguagem.

Não existe um tipo de ponto fixo separado, papel que cabe ao \code{float} configurável, e tampouco existem \code{char}, \code{double}, \code{struct}, \code{union}, \code{enum} ou \textit{strings}; um literal de texto só aparece na inicialização de \textit{arrays} por arquivo, adiante.

\subsection{Literais complexos e o \texttt{i} reservado}

Um literal complexo escreve-se \code{a+bi}, como em matemática:

\begin{lstlisting}[style=cmm]
comp z;
z = 3.0 + 4.0i;
\end{lstlisting}

Por causa disso, o identificador \code{i} é reservado para a unidade imaginária, e usá-lo como nome de variável é erro de compilação. Adote \code{k}, \code{n} ou \code{idx} para contadores.

\subsection{Conversões entre tipos}

Misturar \code{int}, \code{float} e \code{comp} numa expressão ou atribuição é permitido, mas o compilador emite avisos de conversão implícita, por exemplo ao guardar um \code{float} num \code{int} com perda. Usar \code{float} ou \code{comp} como condição de \code{if} ou \code{while}, ou como índice de \textit{array}, também gera aviso. Trate esses avisos como cheiro de projeto: conversões explícitas e índices inteiros deixam o hardware gerado mais barato e o comportamento mais óbvio.

\section{Operadores}

A Tabela~\ref{tab:operadores} reúne os operadores da linguagem, cuja precedência segue a do C.

\begin{table}[ht]
\centering
\caption{Operadores da \cmm{}.}
\label{tab:operadores}
\small
\begin{tabular}{lp{0.56\linewidth}}
\toprule
Grupo & Operadores \\
\midrule
Aritméticos & \texttt{+} \enspace \texttt{-} \enspace \texttt{*} \enspace \texttt{/} \enspace \texttt{\%} (módulo, só inteiros) e o \texttt{-} unário \\
Bit a bit & \texttt{\&} \enspace \texttt{|} \enspace \texttt{\textasciicircum} \enspace \texttt{\textasciitilde} (inversão) \\
Deslocamentos & \texttt{<<} \enspace \texttt{>>} \enspace \texttt{>>>} (à direita, preservando o sinal) \\
Relacionais & \texttt{<} \enspace \texttt{>} \enspace \texttt{<=} \enspace \texttt{>=} \enspace \texttt{==} \enspace \texttt{!=} \\
Lógicos & \texttt{\&\&} \enspace \texttt{||} \enspace \texttt{!} \\
Incremento & \texttt{++} (como comando ou em expressão, inclusive em elemento de \textit{array}) \\
\bottomrule
\end{tabular}
\end{table}

As ausências importam tanto quanto as presenças.

\begin{aviso}
Não existem em \cmm{}: o decremento \texttt{-{}-}, os operadores compostos (\texttt{+=}, \texttt{-=}, \texttt{*=} e afins), o operador ternário \texttt{?:}, o \textit{cast} explícito e os ponteiros; os símbolos \texttt{*} e \texttt{\&} são apenas multiplicação e E bit a bit. Escreva \texttt{x = x - 1;} onde escreveria \texttt{x-{}-;}.
\end{aviso}

\subsection{Operadores custam hardware}

Cada operador usado pela primeira vez no programa instancia o circuito correspondente na ULA do processador gerado, e o terminal TASM anuncia cada instância durante a compilação. Divisão e módulo são os blocos mais caros; os deslocamentos, os mais baratos. Daí o idioma do exemplo condutor, \code{soma >> 2} em vez de \code{soma / 4}. A aritmética de complexos aceita as quatro operações, cada uma expandindo para o conjunto de operações reais e imaginárias necessárias.

\section{Estruturas de controle}

Todas as estruturas clássicas do C existem, com a mesma sintaxe: o \texttt{if} com \texttt{else}; o \texttt{while} e o \texttt{do while}; o \texttt{for}, que o compilador reescreve internamente como um \texttt{while} equivalente; o \texttt{switch} com \texttt{case} e \texttt{default}, incluindo o \textit{fall-through} real do C, em que a execução continua no próximo \texttt{case} na ausência de \code{break}; e \code{break}, \code{continue} e \code{return}, os dois primeiros válidos apenas dentro de laços.

\begin{lstlisting}[style=cmm]
// exemplo: saturador
int satura(int x, int lim)
{
    if (x > lim)  return lim;
    if (x < -lim) return -lim;
    return x;
}
\end{lstlisting}

\section{Funções}

As funções seguem a sintaxe do C, \texttt{tipo nome(tipo p1, tipo p2, \ldots)}, com chamadas por argumentos posicionais e retorno por \code{return}. O número de argumentos é conferido em cada chamada, e a contagem errada é erro de compilação; conversões de tipo em parâmetros e no retorno geram aviso. A profundidade de chamadas aninhadas é limitada pela pilha de sub-rotinas, a diretiva \code{\#SDEPTH}.

Duas restrições merecem destaque. \textit{Arrays} não podem ser passados como parâmetro; trabalhe com \textit{arrays} globais ou no próprio \code{main()}. E a recursão não é suportada no fluxo \cmm{}: cada variável, incluindo parâmetros e locais, vive num endereço fixo da memória de dados, de modo que uma função que chamasse a si mesma sobrescreveria os próprios dados. É justamente esse endereçamento fixo que permite à AURORA mostrar cada variável por nome nas formas de onda. Se o algoritmo é essencialmente recursivo, o caminho C++ atende (Capítulo~\ref{cap:cmm-avancado}).

\begin{lstlisting}[style=cmm]
int add(int a, int b)
{
    return a + b;
}

int triple_sum(int x, int y, int z)
{
    return x + y + z;
}

void main()
{
    while (1)
    {
        int v = in(0);
        out(0, add(v, 1));
        out(0, triple_sum(v, 1, 2));
    }
}
\end{lstlisting}

\section{Arrays}

A linguagem aceita \textit{arrays} de uma e duas dimensões, com tamanho fixo, e uma forma de inicialização peculiar e muito útil:

\begin{lstlisting}[style=cmm]
int  x[128];               // sem inicializacao
int  h[8]    "coefs.txt";  // inicializado por arquivo, em tempo de compilacao
int  m[8][8] "matriz.txt"; // 2D, tambem por arquivo
\end{lstlisting}

A inicialização por arquivo lê os valores de um arquivo de texto ao lado do fonte e os grava diretamente na imagem da memória de dados do processador: os coeficientes do seu filtro já nascem na memória, sem custo de inicialização em tempo de execução.

Há duas formas de indexação. A normal, \code{x[k]}, e a bit-reversa, \code{x[k)}, que se distingue pelo fecha-parênteses: o índice é lido com os bits invertidos, o endereçamento clássico da FFT (\textit{Fast Fourier Transform}), com a quantidade de bits invertidos definida pela diretiva \code{\#FFTSIZ} (Capítulo~\ref{cap:cmm-avancado}). Elementos de \textit{array} aceitam o incremento, como em \code{hist[k]++;}.

\section{Variáveis e escopo}

Declarações podem aparecer no topo do programa ou dentro de funções e blocos. Na prática, porém, todo nome vive num endereço fixo e único da memória de dados, e declarar duas variáveis com o mesmo nome é erro. Pense nas variáveis \cmm{} como registradores nomeados do seu circuito: é exatamente isso que elas se tornam, e é por isso que cada uma aparece por nome na forma de onda da simulação.
